\section{Conclusion}
This paper presented a systematic literature review of large language models and small language models across monolithic, hybrid, multi-agent, ensemble, and retrieval-centered architectures. The review was structured around four questions: comparative performance, efficiency, orchestration, and domain suitability. The reanalysis of the 54-paper website corpus shows a field that is rapidly expanding, strongly concentrated in 2025, and increasingly organized around selective cooperation between small and large models rather than direct head-to-head replacement.

The main technical conclusion is that SLMs are most effective when they are embedded in a well-designed system. Retrieval, routing, uncertainty estimation, task decomposition, prompt optimisation, and multimodal fusion repeatedly appear as the mechanisms that allow small models to recover or approach large-model quality. In many reviewed cases, the decisive variable is not the raw parameter count but the quality of the coordination logic around the model.

The main methodological conclusion is that the literature is ahead in architectural creativity but behind in evaluation standardization. Accuracy is almost universally reported, but latency, cost, and energy are not. For future research, the most urgent needs are standardized resource reporting, better working definitions of what counts as an SLM in practice, and domain-grounded benchmarks that reflect deployment constraints.

From a conference perspective, the strongest message of this review is pragmatic: SLMs are not universal replacements for LLMs, and LLMs are not the inevitable answer to every language task. The reviewed evidence supports a third position in which well-designed hybrid systems deliver the best trade-off between capability, controllability, and deployability.
